% ===============================
% MAT221 -- Real Analysis
% Riemann Integral
% ===============================

\documentclass[12pt,a4paper]{article}

% ===============================
% Metadata
% ===============================
\newcommand*{\CourseCode}{MAT221}
\newcommand*{\CourseTitle}{Real Analysis}
\newcommand*{\Chapter}{Riemann Integral}
\newcommand*{\Topics}{Upper and Lower Sums, Riemann Integrability}
\newcommand*{\DocDate}{August 11, 2026}

% ===============================
% Header
% ===============================
\input{header}

% ===============================
% Document
% ===============================
\begin{document}

\FrontPage
\Large

% ======================================================
\begin{heading}
How Should Integration Be Defined?
\end{heading}

In elementary calculus, integration is often introduced as the inverse
operation of differentiation. For sufficiently nice functions, this
relationship is expressed by the Fundamental Theorem of Calculus.

However, the idea of area should not depend on our ability to find an
antiderivative.

The purpose of this chapter is to define integration directly from
properties of the function itself.

The basic idea will be:

\begin{center}
\emph{Approximate the area using rectangles and make the approximation
arbitrarily accurate.}
\end{center}

\clearpage

% ======================================================
\begin{heading}
The Fundamental Theorem of Calculus
\end{heading}

\vspace{10pt}

\begin{theorem}[Fundamental Theorem of Calculus]
Let $f:[a,b]\to\mathbb{R}$ be continuous. If $F$ is any antiderivative
of $f$ on $[a,b]$, that is,
\[
F'(x)=f(x),
\]
then
\[
\int_a^b f(x)\,dx=F(b)-F(a).
\]
\end{theorem}

\clearpage

\begin{problem}
Let
\[
f(x)=2x+1\]

Use the Fundamental Theorem of Calculus to evaluate
\[
\int_0^2 f(x)\,dx.
\]

First find an antiderivative $F$ of $f$, and then compute
\[
F(2)-F(0).
\]
\end{problem}

\clearpage

% ======================================================
\begin{heading}
Why Antiderivatives Are Not Enough
\end{heading}

Consider the function
\[
h(x)=
\begin{cases}
1, & 0\le x<1,\\
2, & 1\le x\le 2.
\end{cases}
\]

\begin{problem}
Show that $h$ has no antiderivative on $[0,2]$.
\end{problem}

\vfill

Nevertheless, the geometric area under the graph is easy to determine:

\[
1\cdot 1+2\cdot 1=3.
\]

Therefore, it is natural to expect that
\[
\int_0^2 h(x)\,dx=3.
\]

Thus,

\begin{center}
\textbf{integration should be defined independently of
antidifferentiation.}
\end{center}

\clearpage

% ======================================================
\begin{heading}
From Geometry to Approximation
\end{heading}

Suppose that $f$ is a nonnegative function defined on $[a,b]$.

We want to give a precise meaning to the phrase

\begin{center}
\emph{the area under the graph of $f$ from $a$ to $b$.}
\end{center}

The main geometric idea is to divide the interval into smaller pieces
and approximate the region using rectangles.

Some rectangles will underestimate the area, while others will
overestimate it.

If these two types of approximations can be made arbitrarily close,
then there is only one possible value for the area.


\clearpage

% ======================================================
\begin{heading}
Bounded Functions on Closed Intervals
\end{heading}

\begin{definition}[Bounded Function]
A function
\[
f:[a,b]\to\mathbb{R}
\]
is called \textbf{bounded} if there exists $M>0$ such that
\[
|f(x)|\le M
\]
for every $x\in[a,b]$.
\end{definition}

Boundedness is essential in the construction of upper and lower sums.

If a function is unbounded on a subinterval, then a finite upper
rectangular approximation may not exist.

\vspace{15pt}

\begin{theorem}
Every continuous function on a closed interval $[a,b]$ is bounded.
\end{theorem}

\clearpage

% ======================================================
\begin{heading}
Partitions of an Interval
\end{heading}

\begin{definition}[Partition]
A \textbf{partition} $P$ of $[a,b]$ is a finite ordered set
\[
P=\{x_0,x_1,\dots,x_n\}
\]
such that
\[
a=x_0<x_1<\cdots<x_n=b.
\]
\end{definition}


\clearpage

% ======================================================
\begin{heading}
Norm of a Partition
\end{heading}

\begin{definition}[Norm (Mesh)]
Let
\(
P=\{x_0,x_1,\dots,x_n\}
\)
be a partition of $[a,b]$. The \textbf{norm}, or \textbf{mesh}, of $P$ is
\[
|P|
=
\max_{1\le i\le n}
(x_i-x_{i-1}).
\]

\end{definition}

\vspace{100pt}

\begin{problem}
Let
\(
P=
\left\{
0,\frac14,\frac12,1,2
\right\}.
\)

\begin{enumerate}
    \item Compute $|P|$.

    \item Add the point $\frac32$ to obtain the partition
    \(
    Q=
    \left\{
    0,\frac14,\frac12,1,\frac32,2
    \right\}.
    \)

    \item Compute $|Q|$.

    \item Compare $|P|$ and $|Q|$.
\end{enumerate}
\end{problem}

\clearpage

% ======================================================
\begin{heading}
Supremum and Infimum on Subintervals
\end{heading}


\begin{definition}
Let $f$ be bounded on $[a,b]$, and let
\(
P=\{x_0,x_1,\dots,x_n\}
\)
be a partition. For each subinterval $[x_{i-1},x_i]$, define
\[
M_i
=
\sup
\left\{
f(x):x\in[x_{i-1},x_i]
\right\}
\]
and
\[
m_i
=
\inf
\left\{
f(x):x\in[x_{i-1},x_i]
\right\}.
\]
\end{definition}


\clearpage

% ======================================================
\begin{heading}
Lower Sums
\end{heading}

\begin{definition}[Lower Sum]
Let $f$ be bounded on $[a,b]$, and let
\(
P=\{x_0,x_1,\dots,x_n\}
\)
be a partition of $[a,b]$.

The \textbf{lower sum} of $f$ with respect to $P$ is
\[
L(f,P)
=
\sum_{i=1}^{n}
m_i(x_i-x_{i-1}).
\]
\end{definition}



\clearpage

% ======================================================
\begin{heading}
Upper Sums
\end{heading}

\begin{definition}[Upper Sum]
Let $f$ be bounded on $[a,b]$, and let
\[
P=\{x_0,x_1,\dots,x_n\}
\]
be a partition of $[a,b]$.

The \textbf{upper sum} of $f$ with respect to $P$ is
\[
U(f,P)
=
\sum_{i=1}^{n}
M_i(x_i-x_{i-1}).
\]
\end{definition}

\clearpage

\begin{problem}
Let
\[
f(x)=x^{3}-4x^{2}+2
\]
on $[0,4]$, and let
\[
P=\{0,1,2,3,4\}.
\]

\begin{enumerate}
    \item Compute
    \(
    L(f,P).
    \)

    \item Compute
    \(
    U(f,P).
    \)
\end{enumerate}
\end{problem}

\clearpage

% ======================================================
\begin{heading}
Lower Sums Are Below Upper Sums
\end{heading}

\begin{theorem}
For every partition $P$ of $[a,b]$,
\[
L(f,P)\le U(f,P).
\]
\end{theorem}


\clearpage

% ======================================================
\begin{heading}
Refinement of a Partition
\end{heading}

\begin{definition}[Refinement]
Let $P$ and $Q$ be partitions of $[a,b]$. We say that $Q$ is a \textbf{refinement} of $P$ if
\[
P\subseteq Q.
\]
\textbf{Note:} More points in a partition does not necessarily mean that is a refinement. The refinement must contain all the points of the original partition.
\end{definition}


\clearpage

% ======================================================
\begin{heading}
Refinement and Monotonicity
\end{heading}

\begin{theorem}
If $Q$ is a refinement of $P$, then
\[
L(f,P)\le L(f,Q)
\]
and
\[
U(f,Q)\le U(f,P).
\]
\end{theorem}

\clearpage

\begin{problem}
Let
\(
f(x)=x^2
\)
on $[0,4]$.

Consider
\[
P=\{0,2,4\}
\]
and
\[
Q=
\left\{
0,1,2,4
\right\}.
\]

\begin{enumerate}
    \item Verify that $Q$ is a refinement of $P$.

    \item Compute $L(f,P)$ and $U(f,P)$.

    \item Compute $L(f,Q)$ and $U(f,Q)$.

    \item Verify that
    \[
    L(f,P)\le L(f,Q).
    \]

    \item Verify that
    \[
    U(f,Q)\le U(f,P).
    \]

    \item Compare
    \[
    U(f,P)-L(f,P)
    \]
    with
    \[
    U(f,Q)-L(f,Q).
    \]
\end{enumerate}
\end{problem}

\clearpage

% ======================================================
\begin{heading}
Comparing Two Arbitrary Partitions
\end{heading}

\vspace{10pt}

\begin{theorem}
If $P_1$ and $P_2$ are any two partitions (not necessarily refinements) of $[a,b]$, then
\[
L(f,P_1)\le U(f,P_2).
\]
\end{theorem}

Thus, no matter which partitions are chosen, every upper sum is
greater than or equal to every lower sum.

Note: Two partitions do not need to be refinements of each other.

However, given two partitions $P_1$ and $P_2$, their union
\[
P_1\cup P_2
\]
is a partition that refines both of them.

\clearpage
\begin{problem}
Let
\[
f(x)=x^2
\]
on $[0,4]$, and consider
\[
P_1=\{0,1,2,4\}
\]
and
\[
P_2=
\left\{
0,2,3,4
\right\}.
\]

Neither partition is a refinement of the other.

\begin{enumerate}
    \item Compute $L(f,P_1)$.

    \item Compute $U(f,P_2)$.

    \item Verify that
    \[
    L(f,P_1)\le U(f,P_2).
    \]

    \item Find the common refinement
    \[
    P_1\cup P_2.
    \]

    \item Explain the chain
    \[
    L(f,P_1)
    \le
    L(f,P_1\cup P_2)
    \le
    U(f,P_1\cup P_2)
    \le
    U(f,P_2).
    \]
\end{enumerate}
\end{problem}

\clearpage

% ======================================================
\begin{heading}
Upper Integral
\end{heading}

Instead of choosing one particular partition, consider the collection
of upper sums obtained from \emph{all} partitions of $[a,b]$.

Each upper sum is an overestimate.

The smallest possible such overestimate is obtained by taking an
infimum.

\begin{definition}[Upper Integral]
Let $f$ be bounded on $[a,b]$.

The \textbf{upper integral} of $f$ over $[a,b]$ is
\[
\overline{\int_a^b} f
=
\inf
\left\{
U(f,P):
P\text{ is a partition of }[a,b]
\right\}.
\]
\end{definition}

\clearpage

% ======================================================
\begin{heading}
Lower Integral
\end{heading}

Each lower sum is an underestimate.

The largest possible such underestimate is obtained by taking a
supremum.

\begin{definition}[Lower Integral]
Let $f$ be bounded on $[a,b]$.

The \textbf{lower integral} of $f$ over $[a,b]$ is
\[
\underline{\int_a^b} f
=
\sup
\left\{
L(f,P):
P\text{ is a partition of }[a,b]
\right\}.
\]
\end{definition}

\vspace{20pt}

The lower integral is the best possible lower approximation, while
the upper integral is the best possible upper approximation.

\clearpage

% ======================================================
\begin{heading}
Comparing the Lower and Upper Integrals
\end{heading}

\begin{theorem}
For every bounded function
\[
f:[a,b]\to\mathbb{R},
\]
we have
\[
\underline{\int_a^b} f
\le
\overline{\int_a^b} f.
\]
\end{theorem}

\clearpage

\begin{problem}
Let
\[
f(x)=x
\]
on $[0,1]$.

For each $n\in\mathbb{N}$, consider the uniform partition
\[
P_n=
\left\{
0,\frac1n,\frac2n,\dots,\frac{n-1}{n},1
\right\}.
\]

\begin{enumerate}
    \item Show that
    \[
    L(f,P_n)
    =
    \frac{n-1}{2n}.
    \]

    \item Show that
    \[
    U(f,P_n)
    =
    \frac{n+1}{2n}.
    \]

    \item Hence show that
    \[
    U(f,P_n)-L(f,P_n)
    =
    \frac1n.
    \]

    \item Determine
    \[
    \lim_{n\to\infty}
    \left(
    U(f,P_n)-L(f,P_n)
    \right).
    \]

    \item Determine
    \[
    \lim_{n\to\infty}L(f,P_n)
    \qquad\text{and}\qquad
    \lim_{n\to\infty}U(f,P_n).
    \]
\end{enumerate}
\end{problem}

\clearpage

% ======================================================
\begin{heading}
Riemann Integrability
\end{heading}

The previous example suggests the essential idea of integration.

If lower sums and upper sums can be forced arbitrarily close together,
then there should be a unique number representing the integral.

\begin{definition}[Riemann Integrable Function]
A bounded function
\[
f:[a,b]\to\mathbb{R}
\]
is called \textbf{Riemann integrable} on $[a,b]$ if
\[
\underline{\int_a^b} f
=
\overline{\int_a^b} f.
\]

The common value is called the \textbf{Riemann integral} of $f$ and
is denoted by
\[
\int_a^b f(x)\,dx.
\]
\end{definition}

\vspace{20pt}

In other words, $f$ is Riemann integrable when there is no gap between
the best lower approximation and the best upper approximation.


\clearpage


% ======================================================
\begin{heading}
A Pathological Counterexample
\end{heading}

\begin{definition}[Dirichlet Function]
Define
\[
f(x)=
\begin{cases}
1, & x\in\mathbb{Q},\\
0, & x\notin\mathbb{Q}.
\end{cases}
\]
\end{definition}

Both $\mathbb{Q}$ and $\mathbb{R}\setminus\mathbb{Q}$ are dense in
$\mathbb{R}$.

Therefore, every nonempty interval contains both rational and
irrational numbers.

\vspace{15pt}

\begin{problem}
Show that the Dirichlet function is not Riemann integrable on $[0,1]$.
\end{problem}

\clearpage


\EndPage

\end{document}